% Module 8: Quantum Information and Error Correction
% Estimated slides: 30 | Duration: ~2 hours
% Key topics: Classical information theory basics, quantum entropy, no-cloning revisited,
%             error types (bit-flip, phase-flip), basic error correcting codes,
%             surface codes (conceptual), fault tolerance, error mitigation for NISQ
% Labs: Bit-flip error correction circuit in Qiskit

%%%%%%%%%%%%%%%%%%%%%%%%%%%%%%%%%%%%%%%%%%%%%%%%%%%%%%%%%%
\begin{frame}[fragile]\frametitle{}
\begin{center}
{\Large Module 8: Quantum Information and Error Correction}
\end{center}
\end{frame}

%%%%%%%%%%%%%%%%%%%%%%%%%%%%%%%%%%%%%%%%%%%%%%%%%%%%%%%%%%%
\begin{frame}[fragile]\frametitle{Why Quantum Error Correction?}
\begin{itemize}
\item Real qubits are noisy: gate errors, decoherence, measurement errors
\item Error rates: 0.1-1\% per gate operation; accumulate over deep circuits
\item Classical computers: engineers can rely on near-perfect gates (error rate $10^{-15}$)
\item Quantum: errors are much worse --- need active correction during computation
\item Classical error correction: copy bits 3x, take majority vote
\item Quantum challenge: cannot copy qubits (no-cloning) and measuring destroys state
\item Need entirely new error correction paradigm
\end{itemize}
% Suggested diagram: Error accumulation curve vs circuit depth for different error rates
\end{frame}

%%%%%%%%%%%%%%%%%%%%%%%%%%%%%%%%%%%%%%%%%%%%%%%%%%%%%%%%%%%
\begin{frame}[fragile]\frametitle{Types of Quantum Errors}
\begin{itemize}
\item \textbf{Bit-flip error}: $|0\rangle \to |1\rangle$ or $|1\rangle \to |0\rangle$ (like classical bit flip)
\item \textbf{Phase-flip error}: $|+\rangle \to |-\rangle$ (flips the sign of $|1\rangle$ amplitude)
\item \textbf{Depolarizing error}: random mixture of X, Y, Z errors
\item \textbf{Amplitude damping}: qubit loses energy, $|1\rangle$ decays to $|0\rangle$ (T1 error)
\item \textbf{Dephasing}: phase coherence lost over time (T2 error)
\item \textbf{Readout error}: qubit measured incorrectly
\item Any error can be decomposed into X (bit-flip) and Z (phase-flip) components
\end{itemize}
% Suggested diagram: Bloch sphere showing different error types as rotations/decoherence
\end{frame}

%%%%%%%%%%%%%%%%%%%%%%%%%%%%%%%%%%%%%%%%%%%%%%%%%%%%%%%%%%%
\begin{frame}[fragile]\frametitle{Classical Error Correction: Repetition Code}
\begin{itemize}
\item Encode 1 logical bit using 3 physical bits: $0 \to 000$, $1 \to 111$
\item If one bit flips (say $000 \to 010$), majority vote corrects it back to $000$
\item Can detect and correct any 1-bit error
\item Overhead: 3x physical bits per logical bit
\item Key idea: encode information redundantly so errors can be detected and corrected
\item Quantum version needs to handle both bit-flip AND phase-flip errors simultaneously
\end{itemize}
% Suggested diagram: Classical repetition code with error, syndrome, and correction
\end{frame}

%%%%%%%%%%%%%%%%%%%%%%%%%%%%%%%%%%%%%%%%%%%%%%%%%%%%%%%%%%%
\begin{frame}[fragile]\frametitle{Quantum Bit-Flip Code}
\begin{itemize}
\item Encode $|\psi\rangle = \alpha|0\rangle + \beta|1\rangle$ as:
  $\alpha|000\rangle + \beta|111\rangle$ using CNOT gates
\item A bit-flip error on any one of 3 qubits is detectable
\item Measure parity checks: XOR of pairs (01 parity, 12 parity) --- called syndrome measurement
\item Parity checks don't reveal the qubit state, only where the error occurred
\item Apply X correction on the erroneous qubit
\item Corrects any single bit-flip error; phase-flip errors still unprotected
\end{itemize}
% Suggested diagram: 3-qubit encoding circuit; syndrome measurement circuit
\end{frame}

%%%%%%%%%%%%%%%%%%%%%%%%%%%%%%%%%%%%%%%%%%%%%%%%%%%%%%%%%%%
\begin{frame}[fragile]\frametitle{Quantum Phase-Flip Code}
\begin{itemize}
\item Bit-flip code does not protect against phase-flip errors
\item Solution: encode in the X basis ($|+\rangle$ and $|-\rangle$) using H gates
\item Apply H to each qubit before the repetition code; phase flips become bit flips in this basis
\item Shor code (1995): combines bit-flip and phase-flip protection
\item Uses 9 physical qubits to encode 1 logical qubit
\item Corrects any single-qubit error (bit-flip, phase-flip, or both)
\end{itemize}
% Suggested diagram: Shor 9-qubit code encoding structure
\end{frame}

%%%%%%%%%%%%%%%%%%%%%%%%%%%%%%%%%%%%%%%%%%%%%%%%%%%%%%%%%%%
\begin{frame}[fragile]\frametitle{Syndrome Measurement: Detecting Without Collapsing}
\begin{itemize}
\item Key insight: measure parity (even/odd) of pairs, NOT the qubits themselves
\item Parity measurement: $|00\rangle$ and $|11\rangle$ are even; $|01\rangle$ and $|10\rangle$ are odd
\item Parity checks don't reveal individual qubit values --- no state collapse
\item Error syndrome: the pattern of parity check results identifies the error location
\item This is the crucial trick that makes quantum error correction possible
\item Ancilla qubits: used to hold parity measurement results
\end{itemize}
% Suggested diagram: Parity check circuit with ancilla qubits; syndrome table
\end{frame}

%%%%%%%%%%%%%%%%%%%%%%%%%%%%%%%%%%%%%%%%%%%%%%%%%%%%%%%%%%%
\begin{frame}[fragile]\frametitle{Surface Codes: The Leading Approach}
\begin{itemize}
\item Surface codes: qubits arranged in 2D grid; physical qubits connect to nearest neighbors only
\item Data qubits (logical info) + ancilla qubits (for syndrome measurement) alternate in grid
\item Error correction: measure all ancilla qubits, infer error locations from syndrome pattern
\item Threshold theorem: if error rate per gate $< 1\%$ (threshold), can correct to arbitrary precision
\item Overhead: $\sim$1000 physical qubits per 1 logical qubit (for deep circuits)
\item Leading approach for IBM, Google fault-tolerant roadmaps
\end{itemize}
% Suggested diagram: Surface code 2D grid layout with data and ancilla qubits
\end{frame}

%%%%%%%%%%%%%%%%%%%%%%%%%%%%%%%%%%%%%%%%%%%%%%%%%%%%%%%%%%%
\begin{frame}[fragile]\frametitle{The Quantum Error Correction Overhead}
\begin{itemize}
\item Each logical qubit requires many physical qubits (depending on target error rate)
\item Distance-$d$ surface code: $\sim 2d^2$ physical qubits, corrects up to $(d-1)/2$ errors
\item For Shor's algorithm on 2048-bit RSA: $\sim$4000 logical qubits needed
\item At distance 27: $\sim$2916 physical qubits per logical qubit
\item Total: $\sim$12 million physical qubits
\item IBM's roadmap targets 100,000+ physical qubits by 2033; error correction by 2029-2031
\end{itemize}
% Suggested diagram: Physical vs logical qubit ratio curve for different distances
\end{frame}

%%%%%%%%%%%%%%%%%%%%%%%%%%%%%%%%%%%%%%%%%%%%%%%%%%%%%%%%%%%
\begin{frame}[fragile]\frametitle{Error Mitigation: Near-Term Techniques}
\begin{itemize}
\item Error mitigation: reduce noise effects without full error correction
\item Suitable for NISQ devices (no overhead of error correction)
\item \textbf{Zero-noise extrapolation}: run at multiple noise levels, extrapolate to zero noise
\item \textbf{Probabilistic error cancellation}: model noise channels, cancel their effect statistically
\item \textbf{Clifford data regression}: use classical simulation of Clifford circuits as reference
\item \textbf{Readout error mitigation}: measure calibration matrix, invert to correct readout
\item Qiskit: \texttt{qiskit.ignis} and \texttt{mthree} packages for mitigation
\end{itemize}
% Suggested diagram: Zero-noise extrapolation plot showing 3 noise levels and extrapolated value
\end{frame}

%%%%%%%%%%%%%%%%%%%%%%%%%%%%%%%%%%%%%%%%%%%%%%%%%%%%%%%%%%%
\begin{frame}[fragile]\frametitle{Quantum Entropy and Information}
\begin{itemize}
\item Classical Shannon entropy: $H = -\sum p_i \log p_i$ (uncertainty in a distribution)
\item Quantum von Neumann entropy: $S = -\text{Tr}(\rho \log \rho)$; $\rho$ = density matrix
\item Pure state: $S = 0$ (no entropy, perfectly known); maximally mixed: $S = \log d$
\item Entanglement entropy: entropy of subsystem; measures entanglement between parts
\item Bell state: subsystem entropy = 1 bit (maximally entangled)
\item Holevo bound: can extract at most $n$ classical bits from $n$ qubits
\end{itemize}
% Suggested diagram: Entropy visualization: pure vs mixed state density matrices
\end{frame}

%%%%%%%%%%%%%%%%%%%%%%%%%%%%%%%%%%%%%%%%%%%%%%%%%%%%%%%%%%%
\begin{frame}[fragile]\frametitle{Lab: Bit-Flip Error Correction in Qiskit}
\begin{itemize}
\item Encode $|+\rangle$ state using 3-qubit repetition code
\item Introduce a controlled bit-flip error on qubit 1
\item Perform syndrome measurement using 2 ancilla qubits
\item Apply correction based on syndrome; verify state is recovered
\end{itemize}
\begin{verbatim}
# Encode: CNOT(0->1), CNOT(0->2)
# Error: qc.x(1)  # flip qubit 1
# Syndrome: CNOT(0->anc0) XOR CNOT(1->anc0); etc.
# Correct: apply X to qubit 1 if syndrome = 10
\end{verbatim}
\end{frame}

%%%%%%%%%%%%%%%%%%%%%%%%%%%%%%%%%%%%%%%%%%%%%%%%%%%%%%%%%%%
\begin{frame}[fragile]\frametitle{Module 8 Summary}
\begin{itemize}
\item Quantum errors: bit-flip (X), phase-flip (Z), depolarizing, amplitude damping
\item Classical error correction ideas don't directly apply (no-cloning blocks copying)
\item Syndrome measurement: detect errors without collapsing the logical state
\item 3-qubit code corrects bit-flip errors; Shor's 9-qubit code corrects all single-qubit errors
\item Surface codes: leading approach for fault-tolerant QC; $\sim$1000 physical per logical qubit
\item Error mitigation: practical near-term techniques for NISQ devices
\item Threshold theorem: below ~1\% gate error rate, error-free computation is possible
\end{itemize}
\end{frame}
